\clearpage
\newcounter{ch2checkpointappendix}
\renewcommand{\thech2checkpointappendix}{\thechapter.A}
\refstepcounter{ch2checkpointappendix}
\section*{Chapter 2 Appendix: Checkpoint-Selection Diagnostics}
\addcontentsline{toc}{section}{Chapter 2 Appendix: Checkpoint-Selection Diagnostics}
\label{app:checkpoint_selection}

This appendix records generation-stability and task-aligned checkpoint
diagnostics for \textit{Model chk-1}, \textit{Model chk-2}, and \textit{Model
chk-3}. The chk-1 gate
uses a frozen eight-case generation probe; checkpoints 200 and 255 were tested
contemporaneously, whereas the intermediate post-cp200 checkpoints were replayed
retrospectively. The chk-2 gate uses a deterministic twelve-observation sample
that was held out from gradient updates but participated directly in checkpoint
selection. Neither diagnostic is an independent population test, and neither
supports population-level inference. The chk-2 reference paragraphs are
synthetic teacher targets rather than official FOMC Minutes.

The appendix reports the checkpoint gates in separate tables because the model
stages use different task contracts. Table~\ref{tab:app:chk1_post_cp200_hard_gates} audits
free-running analysis generation, whereas
Table~\ref{tab:app:chk2_exact_merged_n12_selection} evaluates exact-merged
Minutes rewriting. Tables~\ref{tab:app:chk3_sft_qualification_gates}
and~\ref{tab:app:chk3_grpo_selection_gates} distinguish the decision-SFT
qualification gate from the subsequent GRPO checkpoint-selection gate, while
Table~\ref{tab:app:chk3_sft_cp38_cp39_selection} records why SFT checkpoint 38,
rather than the terminal checkpoint 39, initializes GRPO. Their case counts and
pass rates are therefore not directly comparable.

\subsection*{Learning-Rate Stability Diagnostics for Model chk-1}
\label{app:chk1_lr_selection}

The peak learning rate for \textit{Model chk-1} was chosen through staged
generation-stability diagnostics rather than an exhaustive hyperparameter
search. The first controlled smoke comparison held the dataset, sample order,
random seeds, LoRA configuration, optimizer, warm-up schedule, batch size, and
generation prompts fixed while reducing the peak learning rate from
\(5\times10^{-5}\) to \(1\times10^{-5}\). At \(5\times10^{-5}\), the first
matched held-out completion reached the 3,072-token limit, failed to emit both
an end-of-sequence token and a valid reasoning boundary, and recorded full-text
and tail four-gram repetition rates of 0.9231 and 0.9647, respectively. This
single case satisfied the preregistered fail-fast condition, so the remaining
seven cases were not generated. By contrast, the ten-step
\(1\times10^{-5}\) smoke checkpoint passed all eight fixed probe cases.

Passing the short smoke test did not establish stability over a complete
training trajectory. The three-epoch \(1\times10^{-5}\) run completed all 255
optimizer updates and attained a terminal validation loss of 1.300019, but its
terminal checkpoint failed the same generation gate: one of eight completions
reached the token limit, only seven satisfied the output contract, and three
were classified as catastrophically repetitive. The lower validation loss
therefore did not imply reliable free-running generation. The peak learning
rate was subsequently reduced to \(1\times10^{-6}\), and a new three-epoch run
was initialized from \textit{Model chk-0}. In that run, checkpoints 80, 150,
200, and 255 all passed the fixed eight-case generation probe. Checkpoint 200
was retained as the paper-facing \textit{Model chk-1} because it combined the
minimum logged validation loss (1.686099) with a fully passing generation
probe; checkpoint 255 had an essentially unchanged terminal validation loss
of 1.686131 and also passed, but exhibited slightly greater mean four-gram
repetition (0.140996 versus 0.134664).

\begin{table}[htbp]
    \centering
    \scriptsize
    \setlength{\tabcolsep}{3.5pt}
    \renewcommand{\arraystretch}{1.15}
    \begin{tabularx}{\textwidth}{@{}l>{\raggedright\arraybackslash}Xcccccc@{}}
        \toprule
        \makecell{\textbf{Peak learning}\\\textbf{rate}} &
        \textbf{Diagnostic state} &
        \textbf{Cases} &
        \makecell{\textbf{Contract}\\\textbf{valid}} &
        \textbf{Capped} &
        \makecell{\textbf{Catastrophic}\\\textbf{repetition}} &
        \makecell{\textbf{Mean full-text}\\\textbf{4-gram repetition}} &
        \textbf{Gate} \\
        \midrule
        \(5\times10^{-5}\) & ten-step smoke, cp10 & 1\textsuperscript{a} & 0/1 & 1/1 & 1/1 & 0.923102 & Failed \\
        \(1\times10^{-5}\) & ten-step smoke, cp10 & 8 & 8/8 & 0/8 & 0/8 & 0.094759 & Passed \\
        \(1\times10^{-5}\) & three-epoch run, cp255 & 8 & 7/8 & 1/8 & 3/8 & 0.408608 & Failed \\
        \(1\times10^{-6}\) & three-epoch run, cp200 & 8 & 8/8 & 0/8 & 0/8 & 0.134664 & Passed \\
        \bottomrule
    \end{tabularx}
    \begin{minipage}{0.96\textwidth}
        \vspace{2pt}
        \footnotesize
        \textit{Note:} \textsuperscript{a}\ The \(5\times10^{-5}\) probe stopped
        after its first case triggered the fail-fast rule; its row is therefore
        a failure demonstration, not an estimate of an eight-case failure rate.
        ``Catastrophic repetition'' follows the frozen probe thresholds of at
        least 0.50 full-text or 0.60 tail token four-gram repetition. All rows
        use the same frozen probe manifest, but they represent an adaptive,
        staged diagnostic sequence rather than a balanced or statistically
        powered learning-rate sweep.
    \end{minipage}
    \caption{Staged learning-rate stability diagnostics for \textit{Model
    chk-1}. Selection prioritized valid and non-degenerate free-running
    generation; validation loss alone was not treated as sufficient evidence
    of model reliability.}
    \label{tab:app:chk1_lr_stability}
\end{table}

\subsection*{Post-cp200 Generation-Stability Replay for Model chk-1}

To determine whether the retained checkpoint was uniquely stable or whether
stability persisted across the late loss plateau, checkpoints 210--250 were
replayed retrospectively on 23 August 2026 using the same frozen eight-case
manifest, seeds, decoding settings, \textit{Model chk-0} baseline, and runtime
versions as the original checkpoint-200 and checkpoint-255 probes. These
intermediate replays did not participate in the original checkpoint-selection
decision. Table~\ref{tab:app:chk1_post_cp200_hard_gates} therefore provides a
matched descriptive audit of the late trajectory rather than new prospective
selection evidence.

\begin{table}[htbp]
    \centering
    \scriptsize
    \setlength{\tabcolsep}{3.25pt}
    \renewcommand{\arraystretch}{1.15}
    \begin{tabular}{@{}lcccccc@{}}
        \toprule
        \textbf{Checkpoint} &
        \makecell{\textbf{Training / evaluation}\\\textbf{loss}} &
        \makecell{\textbf{Finite / EOS /}\\\textbf{contract-valid}} &
        \makecell{\textbf{Cap / catastrophic /}\\\textbf{periodic}} &
        \makecell{\textbf{Mean full-text}\\\textbf{4-gram repetition}} &
        \makecell{\textbf{$\Delta$ versus chk-0}\\\textbf{mean / maximum}} &
        \textbf{Gate} \\
        \midrule
        \texttt{cp200} (retained) & 1.729560 / 1.686099 & 8 / 8 / 8 & 0 / 0 / 0 & 0.134664 & 0.051945 / 0.144158 & Passed \\
        \texttt{cp210}\textsuperscript{a} & 1.694860 / 1.686223 & 8 / 8 / 8 & 0 / 0 / 0 & 0.148501 & 0.065781 / 0.135519 & Passed \\
        \texttt{cp220}\textsuperscript{a} & 1.709530 / 1.686188 & 8 / 8 / 8 & 0 / 0 / 0 & 0.142607 & 0.059888 / 0.161819 & Passed \\
        \texttt{cp230}\textsuperscript{a} & 1.724720 / 1.686177 & 8 / 8 / 8 & 0 / 0 / 0 & 0.116562 & 0.033842 / 0.158795 & Passed \\
        \texttt{cp240}\textsuperscript{a} & 1.704090 / 1.686112 & 8 / 8 / 8 & 0 / 0 / 0 & 0.111440 & 0.028720 / 0.146266 & Passed \\
        \texttt{cp250}\textsuperscript{a} & 1.705740 / 1.686216 & 8 / 8 / 8 & 0 / 0 / 0 & 0.150136 & 0.067417 / 0.141370 & Passed \\
        \texttt{cp255} (terminal) & 1.722860 / 1.686131 & 8 / 8 / 8 & 0 / 0 / 0 & 0.140996 & 0.058276 / 0.171121 & Passed \\
        \bottomrule
    \end{tabular}
    \begin{minipage}{0.97\textwidth}
        \vspace{2pt}
        \footnotesize
        \textit{Note:} \textsuperscript{a}\ Retrospective matched replay; the
        original selection record contained probes for checkpoints 200 and 255
        but not for these intermediate checkpoints. Each three-part count is
        the number of cases, out of eight, in the order stated in its column
        heading. Each probe contains six sampled and two greedy cases. The
        frozen comparator requires finite
        output for every case, no completion at the 3,072-token cap, no strict
        periodic tail, no catastrophic repetition (full-text four-gram rate
        $\geq 0.50$ or tail rate $\geq 0.60$), no contract-validity loss relative
        to the fully valid chk-0 baseline, a mean repetition increase no greater
        than 0.10, and a per-case increase no greater than 0.20. EOS is reported
        as a diagnostic but is not a separate comparator threshold. In the loss
        column, the training value is the trailing mean of the ten step-level
        losses ending at the reported checkpoint,
        $\overline{L}^{\mathrm{train}}_{c,10}=10^{-1}\sum_{j=0}^{9}
        L^{\mathrm{train}}_{c-j}$, followed by evaluation loss. Neither loss is
        part of the generation gate; the checkpoint-255 evaluation value is the
        terminal evaluation loss. Full-text and 1,024-token-tail repetition
        values coincide in every row because no retained completion contains
        more than 1,023 content tokens.
    \end{minipage}
    \caption{Post-cp200 generation-stability hard-gate comparison for
    \textit{Model chk-1}. All evaluated checkpoints pass the frozen gate, so
    the table does not establish checkpoint 200 as uniquely stable. Its
    retention remains supported by the minimum logged evaluation loss together
    with a passing contemporaneous generation probe.}
    \label{tab:app:chk1_post_cp200_hard_gates}
\end{table}

\subsection*{Hard-Gate Candidate Selection for Model chk-2}

Table~\ref{tab:app:chk2_exact_merged_n12_selection} reports the deployment-form
replay of the exact-merged candidates from the repository \texttt{chk3}
Minutes-rewriting trajectory. Repository \texttt{chk3} checkpoint 318 is
denoted \textit{Model chk-2} in the paper. Adapter screening and exact-merged
replay are kept distinct because merge form and quantized loading can affect
boundary cases.

\begin{table}[htbp]
    \centering
    \scriptsize
    \setlength{\tabcolsep}{3pt}
    \renewcommand{\arraystretch}{1.15}
    \begin{tabular}{@{}lccccccc@{}}
        \toprule
        \textbf{Candidate} &
        \makecell{\textbf{Training / evaluation}\\\textbf{loss}} &
        \makecell{\textbf{Core}\\\textbf{hard-valid}} &
        \textbf{EOS} &
        \makecell{\textbf{Cap /}\\\textbf{periodic}} &
        \makecell{\textbf{Numeric /}\\\textbf{date valid}} &
        \makecell{\textbf{4-gram repetition}\\\textbf{mean full / max tail}} &
        \textbf{Gate} \\
        \midrule
        \texttt{cp200} & \makecell{1.163500 \\ 1.122187} & 11/12 & 12/12 & 0 / 0 & 12/12 / 11/12 & \makecell{0.150280 \\ 0.346863} & Failed \\
        \makecell[l]{\texttt{cp250} \\(minimum eval. loss)} & \makecell{1.171210 \\ 1.121204} & 10/12 & 11/12 & 1 / 1 & 11/12 / 10/12 & \makecell{0.267695 \\ 0.947111} & Failed \\
        \makecell[l]{\texttt{cp318} \\(retained)} & \makecell{1.158010 \\ 1.121280} & 12/12 & 12/12 & 0 / 0 & 12/12 / 12/12 & \makecell{0.144111 \\ 0.335354} & Passed \\
        \bottomrule
    \end{tabular}
    \begin{minipage}{0.97\textwidth}
        \vspace{2pt}
        \footnotesize
        \textit{Note:} Values in the loss column are the ten-step trailing-mean
        training loss followed by evaluation loss; checkpoint 318 uses the
        terminal evaluation result. Loss is not a hard gate. The probe contains
        twelve deterministic cases. Core hard-valid requires structure,
        numerical fidelity, date fidelity, and non-degeneration to pass jointly
        for a case; a candidate passes the checkpoint gate only when all twelve
        cases are core hard-valid. Adapter screening and exact-merged replay are
        distinct evaluation forms; this table reports only the exact-merged
        deployment-form replay. ``Retained'' denotes the paper-facing evaluated
        checkpoint; the selection receipt is evaluation-only and does not
        authorize canonical promotion or downstream training.
    \end{minipage}
    \caption{Exact-merged checkpoint-selection hard gates for \textit{Model
    chk-2} (repository artifact \texttt{chk3}). Checkpoint 318 is retained
    because it is the only finalist that passes all twelve task-aligned cases,
    despite checkpoint 250 having a marginally lower evaluation loss.}
    \label{tab:app:chk2_exact_merged_n12_selection}
\end{table}

Checkpoint 250 attains the minimum validation loss, but it contains one capped
periodic completion and additional numerical- and date-fidelity failures.
Checkpoint 318 has a validation loss only 0.00007641 above that minimum and is the
only exact-merged finalist to pass all twelve core cases. The evidence therefore
supports checkpoint 318 as a hard-gate-qualified representative of the late
training plateau; it does not establish checkpoint 318 as a statistically
unique optimum.

\subsection*{Model-Level Task-Aligned Diagnostic}

Table~\ref{tab:app:minutes_native_n12_diagnostic} reproduces the model-level
diagnostic used alongside the candidate replay. The experimental comparison is
between the \textit{Model chk-1} checkpoint-200 baseline and \textit{Model
chk-2}, which is repository \texttt{chk3} checkpoint 318. \textit{Model chk-0}
is included only as an auxiliary reference.

\begin{table}[htbp]
    \centering
    \footnotesize
    \setlength{\tabcolsep}{4pt}
    \renewcommand{\arraystretch}{1.15}
    \begin{tabular}{lcccccc}
        \toprule
        \textbf{Model} & \makecell{Structure\\valid} &
        \makecell{Numeric\\valid} & \makecell{Date\\valid} &
        \makecell{Non-degenerate} & \makecell{BERTScore\\F1} &
        \makecell{MPNet\\cosine} \\
        \midrule
        \makecell[l]{\textit{Model chk-0}\\auxiliary reference} &
        11/12 & 9/12 & 10/12 & 11/12 & 0.729516 & 0.721363 \\
        \makecell[l]{\textit{Model chk-1}\\checkpoint 200 (baseline)} &
        12/12 & 12/12 & 12/12 & 12/12 & 0.966947 & 0.944789 \\
        \makecell[l]{\textit{Model chk-2}\\repository \texttt{chk3}, checkpoint 318} &
        12/12 & 12/12 & 12/12 & 12/12 & \textbf{0.973533} & \textbf{0.971603} \\
        \bottomrule
    \end{tabular}
    \caption{Model-level analysis-to-Minutes checkpoint-selection diagnostic
    (\(N=12\)). Semantic similarity is set to zero when an observation fails
    at least one validity criterion. The reference texts are synthetic teacher
    targets rather than official FOMC Minutes.}
    \label{tab:app:minutes_native_n12_diagnostic}
\end{table}

One source analysis is identical to its normalized synthetic target. After
excluding that identity case, the non-identity subsample contains \(N=11\)
observations. \textit{Model chk-2} records BERTScore F1 of 0.975129 and MPNet
cosine similarity of 0.972733, compared with 0.968356 and 0.944044 for
\textit{Model chk-1}. The corresponding paired mean differences are 0.006773
and 0.028688. \textit{Model chk-2} has the higher BERTScore in 7 of 11
comparisons and the higher MPNet similarity in 9 of 11.

These semantic values are descriptive checkpoint-selection diagnostics. The
same deterministic twelve-observation sample was inspected when checkpoint 318
was retained, so neither the \(N=12\) hard-gate table nor the \(N=11\) semantic
comparison is an independent evaluation or a basis for population-level
significance claims.

\clearpage
\subsection*{Separate SFT and GRPO Hard Gates for Model chk-3}

Paper-level \textit{Model chk-3} corresponds to repository artifact
\texttt{chk4}. Its two training stages answer different gate questions. The
decision-SFT gate asks whether an SFT candidate is sufficiently well behaved to
initialize GRPO; the matched replay below distinguishes checkpoint 38 from the
terminal checkpoint 39. The GRPO gate asks whether any later policy checkpoint
meets the frozen checkpoint-selection requirements. Passing the first gate does
not imply passing the second.

\begin{table}[htbp]
    \centering
    \scriptsize
    \setlength{\tabcolsep}{3.5pt}
    \renewcommand{\arraystretch}{1.15}
    \begin{tabularx}{\textwidth}{@{}>{\raggedright\arraybackslash}Xccc>{\raggedright\arraybackslash}X@{}}
        \toprule
        \textbf{Gate component} &
        \textbf{Frozen requirement} &
        \makecell{\textbf{Observed at}\\\textbf{SFT cp38}} &
        \textbf{Result} &
        \textbf{Function of the gate} \\
        \midrule
        Reasoning-boundary delivery & Rate $\geq 0.75$ & 14/15 = 0.9333 & Passed & Requires a usable native reasoning boundary before the terminal action. \\
        Completion cap & Rate $\leq 0.25$ & 1/15 = 0.0667 & Passed & Limits truncated generations under the 1,024-token probe ceiling. \\
        Strict periodic tail & Count $=0$ & 0/15 & Passed & Rejects mechanically looping completions. \\
        Sampled direction correctness by class & At least one per class & \makecell{hold 3; hike 1; cut 1} & Passed & Requires non-zero action signal for each policy direction. \\
        Sampled non-zero reward by class & At least one per class & \makecell{hold 3; hike 1; cut 2} & Passed & Prevents an all-zero reward group. \\
        Sampled reward dispersion by class & Standard deviation $>0$ & \makecell{hold 0.4330; hike 0.1028;\\cut 0.3074} & Passed & Ensures within-group reward variation for GRPO. \\
        \bottomrule
    \end{tabularx}
    \begin{minipage}{0.97\textwidth}
        \vspace{2pt}
        \footnotesize
        \textit{Note:} The probe contains three direction-stratified groups,
        each with one greedy and four sampled generations. Class-specific gates
        use the four sampled generations. Exact magnitude is reported elsewhere
        as a diagnostic but is not an SFT qualification gate, because the
        validation-only 50/75-basis-point hike magnitudes are absent from the
        sealed training split. Checkpoint 38 passes this qualification gate and
        is exact-merged as the pre-GRPO SFT state of \textit{Model chk-3}.
    \end{minipage}
    \caption{Decision-SFT qualification hard gates for \textit{Model chk-3}
    at checkpoint 38. These gates determine whether the SFT state may initialize
    GRPO; they do not establish held-out decision quality.}
    \label{tab:app:chk3_sft_qualification_gates}
\end{table}

\subsubsection*{Why Checkpoint 38 Rather than Checkpoint 39?}

The terminal SFT checkpoint is not automatically preferred to an earlier
checkpoint. Checkpoints 38 and 39 were evaluated with the same
direction-stratified qualification probe and the same frozen thresholds.
Table~\ref{tab:app:chk3_sft_cp38_cp39_selection} separates the trainer-reported
loss observations from the behavioral quantities that actually determine
qualification.

\begin{table}[htbp]
    \centering
    \scriptsize
    \setlength{\tabcolsep}{3.5pt}
    \renewcommand{\arraystretch}{1.15}
    \begin{tabularx}{\textwidth}{
        @{}
        >{\raggedright\arraybackslash}X
        >{\raggedright\arraybackslash}X
        >{\centering\arraybackslash}p{0.22\textwidth}
        >{\centering\arraybackslash}p{0.22\textwidth}
        @{}
    }
        \toprule
        \textbf{Diagnostic} &
        \textbf{Frozen requirement} &
        \textbf{SFT cp38} &
        \textbf{SFT cp39} \\
        \midrule
        Per-step training loss
        & Not a qualification gate
        & 1.6149
        & 1.3127 \\
        Scheduled evaluation loss
        & Not a qualification gate
        & Not recorded at step 38
        & 1.601505 \\
        Reasoning-boundary delivery
        & Overall rate $\geq 0.75$
        & $14/15=0.9333$
        & $13/15=0.8667$ \\
        Completion cap
        & Overall rate $\leq 0.25$
        & $1/15=0.0667$
        & $2/15=0.1333$ \\
        Strict periodic tail
        & Count $=0$
        & 0
        & 0 \\
        \makecell[l]{Sampled direction-correct\\coverage: hold / hike / cut}
        & At least one per class
        & $3 / 1 / 1$
        & $2 / 1 / \mathbf{0}$ \\
        \makecell[l]{Sampled non-zero-reward\\coverage: hold / hike / cut}
        & At least one per class
        & $3 / 1 / 2$
        & $3 / 2 / 3$ \\
        \makecell[l]{Sampled reward standard deviation:\\hold / hike / cut}
        & Greater than zero for every class
        & $0.4330 / 0.1028 / 0.3074$
        & $0.4878 / 0.4263 / 0.0217$ \\
        Overall qualification
        & Every behavioral gate must pass
        & \textbf{Passed}
        & \makecell{\textbf{Failed}: no sampled\\cut-direction success} \\
        \bottomrule
    \end{tabularx}
    \begin{minipage}{0.97\textwidth}
        \vspace{2pt}
        \footnotesize
        \textit{Note:} Each checkpoint is evaluated on one hold, one hike, and
        one cut training prompt, with one greedy and four sampled generations
        per prompt. The class-specific gates use the four sampled generations.
        Checkpoint 39 produces a correct cut under greedy decoding but produces
        no correct cut direction in its four sampled generations. Training loss
        is a batch-level statistic, and no scheduled evaluation was performed at
        step 38; neither loss quantity is part of this behavioral gate.
    \end{minipage}
    \caption{Matched SFT qualification comparison between checkpoints 38 and
    39 for \textit{Model chk-3}. Checkpoint 39 is the terminal SFT state but
    fails the frozen sampled cut-direction coverage requirement.}
    \label{tab:app:chk3_sft_cp38_cp39_selection}
\end{table}

Although checkpoint 39 records the minimum logged SFT evaluation loss, that
value is not a matched checkpoint-38 comparison because evaluation was
scheduled only at steps 13, 26, and 39. More importantly, loss is not the
selection criterion for GRPO initialization. Checkpoint 39 fails the frozen
requirement of at least one sampled direction-correct output in every policy
class: none of its four sampled cut completions predicts the correct direction.
Checkpoint 38 satisfies every qualification component and is therefore retained
and exact-merged as the pre-GRPO state of \textit{Model chk-3}. This small,
train-stratified probe establishes readiness to initialize GRPO, not population
superiority or a statistically unique loss optimum.

\begin{table}[htbp]
    \centering
    \scriptsize
    \setlength{\tabcolsep}{3.5pt}
    \renewcommand{\arraystretch}{1.15}
    \begin{tabular}{@{}lccccccc@{}}
        \toprule
        \textbf{GRPO checkpoint} &
        \makecell{\textbf{Mean}\\\textbf{reward}} &
        \makecell{\textbf{Direction}\\\textbf{accuracy}} &
        \makecell{\textbf{Exact-action}\\\textbf{accuracy}} &
        \makecell{\textbf{Delivery-valid}\\\textbf{rate}} &
        \textbf{Cap rate} &
        \makecell{\textbf{Periodic-tail}\\\textbf{count}} &
        \textbf{Gate} \\
        \midrule
        \texttt{cp310} & 0.227083 & 0.277778 & 0.222222 & 0.444444 & 0.111111 & 0 & Failed \\
        \texttt{cp360} & 0.106944 & 0.222222 & 0.166667 & 0.388889 & 0.000000 & 0 & Failed \\
        \texttt{cp410} & 0.093750 & 0.111111 & 0.111111 & 0.555556 & 0.055556 & 0 & Failed \\
        \texttt{cp440}\textsuperscript{a} & 0.122917 & 0.166667 & 0.111111 & 0.388889 & 0.166667 & 0 & Failed \\
        \texttt{cp450}\textsuperscript{b} & 0.143750 & 0.222222 & 0.166667 & 0.333333 & 0.166667 & 0 & Failed \\
        \texttt{cp460}\textsuperscript{a} & 0.078819 & 0.111111 & 0.055556 & 0.333333 & 0.111111 & 0 & Failed \\
        \texttt{cp468} (terminal) & 0.069444 & 0.055556 & 0.055556 & 0.333333 & 0.166667 & 0 & Failed \\
        \bottomrule
    \end{tabular}
    \begin{minipage}{0.97\textwidth}
        \vspace{2pt}
        \footnotesize
        \textit{Note:} The selection-panel eligibility gate requires a
        delivery-valid rate of at least 0.75, a cap rate no greater than 0.25,
        and zero periodic-tail cases. The reported panel contains 18 completions
        from nine independent-validation prompts; it contains no cut example.
        Every checkpoint fails the delivery threshold. \textsuperscript{a}
        Checkpoints 440 and 460 are post-hoc local-neighbour probes.
        \textsuperscript{b}\ Checkpoint 450 is the best-observed exploratory
        candidate, not a selected checkpoint. On its blind follow-up it records
        reward 0.1725, direction accuracy 0.30, exact-action accuracy 0.15,
        delivery validity 0.30, cap rate 0.10, and zero periodic tails. It covers
        cut, hold, and hike on the train-only retention panel, but fails the
        blind delivery gate. The authoritative receipt therefore records
        \texttt{best\_observed\_checkpoint\_step=450},
        \texttt{selected\_checkpoint\_step=null}, and
        \texttt{provisional\_no\_candidate\_passed\_all\_gates}. The later
        13-meeting test likewise records \texttt{quality\_not\_demonstrated}.
    \end{minipage}
    \caption{GRPO checkpoint-selection hard gates for \textit{Model chk-3}.
    No checkpoint satisfies the frozen eligibility rule; checkpoint 450 was
    evaluated only as the best-observed exploratory candidate through an
    explicit override.}
    \label{tab:app:chk3_grpo_selection_gates}
\end{table}
